\subsection{Second Letter}

\begin{quotex}
In this letter, \textbf{Rene Guenon} chit chats about some of their reading material. He also quotes a long passage from one of \textbf{Julius Evola}'s essays in which Evola again accused Guenon of rationalism, at least by his own definition. Evola insists the principle of the absolute is Power, or Shakti. It is curious that Evola exalts the feminine element above the masculine. Guenon seems annoyed about what he regards as a misunderstanding. Evola also distances himself from Guenon by praising science as the consequence of self-assertion. 

\end{quotex}
Paris, 26 January 1926

Dear Sir,

I believe I am getting later and later with my correspondence; I would have rather liked to give you my wishes for better health since, after you told me in your last letter, you were still suffering.

I am happy that you were satisfied with reading about \textbf{Milarepa}; besides, the contrary would have certainly astonished me. In the same collection, there is another of Bacot's transations: \emph{Three mysterious Tibetans}: that is also good, but these are texts of a somewhat literary character, therefore of much less interest for us.

I suppose that P. Huc's book that you told me about is the republication which was recently made about his travels in Tibet and China; it contains nothing about doctrines; there are in it some curious descriptions, but that is all.

Boehme's work which was translated into Italian must be \emph{De signature rerum}: it seems to me that the meaning of the title was not made very exactly. As to Campanella's book, I am not familiar with it; I have only read a review in the last issues of ``Ignis". The fragments of the Zohar published by Rieder are in effect only extracts from Jean de Pauly's translation; I haven't seen them, but I hear that they were rather well chosen.

As to Mead, his books are so impregnated with Theosophy that I strongly doubt that anything good can be pulled from them.

As for Theosophy, do you know that the solemn proclamation of the new Messiah (Krishnamurti) has to take place very soon with an extraordinary mise en scene? None of that can be taken seriously.

You ask me about the journal ``Ultra"; it defines itself as ``independent theosophical", that is, it is the organ of a group which has separated from Annie Besant's Theosophical Society. Here is the copy of Evola's note which appeared in that review (in an article titled \emph{The Problem of the East and West}), which is about me.

\begin{quotex}
It is clear here that we speak of the East and the West as two ideal types which, if they belong to the general character of the two cultures, cannot belong to their details. Presently, for example, currents like intuitionalism, actual idealism, irrationalism are as a general rule justifiable by the principle of \emph{eros}. We must on the contrary have reservations about pragmatism. Consequently, we can't agree with the thesis held by Rene Guenon (East and West), that ``scientism" and sensual life are the hidden elements connecting them together.

It is correct that western science is used for practical ends. But it is not necessary to confuse the two things: without simple desires and needs which would have remained such, scientific comprehension of nature would not have advanced one step. The accomplishment of science is linked in fact, as we have already noted, to a first manifestation of the positive principle of domination and assertion of the Self—which constitutes a \emph{metaphysical value}. What was made possible by such a principle superior in itself both to the life of needs and to the sentiment that `pure intellectuality' which Guenon, as the good rationalist that he is, has as a superstition—whether it was utilized by elements which are alien to it, is a totally different question. In itself it remains a value that confers on western culture its significance and originality. We assert with Guenon that the principle of the absolute has nothing to do with the sentimental, moralist, and abstractly rational element; but we assert against Guenon that it has nothing to do with that ``pure intellectually" or `metaphysic' that he speaks about, and, emphatically, we \emph{know} what it is. It is for that reason precisely that we dispute, from a higher point of view, that such `metaphysic' can be distinguished from the rational as it was defined by Hegel, for example Venunft opposed to Verstand. And we assert therefore that the \emph{principle of the absolute is power} (\textbf{Shakti}) and whatever system that in the metaphysical order puts something before or above power is rationalistic (in the pejorative sense used by Guenon) and abstract. 

\end{quotex}
Evola does not lack any pretentions as you can see; but, for my part, I still think that he doesn't understand at all what we mean by ``intellectuality", ``knowledge", ``contemplation", etc., and he can't even make the distinction between the ``initiatic" and ``profane" points of view. It appears that he intends to publish a review of my work on the Vedanta in the journal \emph{L'Idealismo Realistico}; we will see what that will be like. In any case, in spite of everything that we tried to explain to him, he persists in finding ``rationalism" in the Vedanta, all while being compelled to recognize that he then takes the word ``rationalism" in a rather different sense from what is usually meant by it.

You must have seen that the publication of \emph{Ignis} was interrupted; Reghini could no longer manage to do it all on his own or pretty much so, or it was clearly necessary that he had to be concerned only with that, which is not possible for him with the conditions of current life. There are certainly people who had promised to help him, but who have done nothing; it is often that way, unfortunately. I think that you would have seen, in the last issue, my article on \textbf{Joseph de Maistre}, concerning the books of Emile Dermenghem\footnote{\url{https://en.wikipedia.org/wiki/Aissawa}}.

I am going to send you the issue of \emph{Candide} containing the article about me which I told you about in my preceding letter. Apart from that, I don't have much of anything new to tell you; the reviews in the journals are a long time coming; and then it must be that people feel some difficulty speaking of my last book. When I have something which will be worth the trouble, I will tell you.

What you say about the conception of Dharma among the Jains seems to me correct, and I don't believe anymore that we can find an essential difference with orthodox doctrine in it. On the contrary, the meaning they give the words \emph{loka} and \emph{aloka} is rather particular; it seems to be quite exact to translate the first, in this case, by the ``world of forms" as you do.

As to \emph{sphota}\footnote{\url{https://en.wikipedia.org/wiki/Spho\%E1\%B9\%ADa}}, it is a conception typical of certain grammarians, and whose discussion which you allude to shows its uselessness. Naturally, the point of view of the grammarians can only be exterior and analytic; it relates uniquely to the form of the word, in as much as this is constituted by the assemblage of certain phonetic elements; the word is not taken there in itself and synthetically. It is well that there is something artificial in the point of view which necessitates the intervention of the \emph{sphota}; you are therefore completely correct. It is also certain that this discussion has much more importance than one could believe when envisaged superficially; but it is not to the current philologists that we can hope make that understood; everything that concerns the true nature of the language eludes them entirely.

Last month I lectured on Oriental metaphysics\footnote{\url{https://www.gornahoor.net/library/OrientalMetaphysics.pdf}} to a study group which meets at the Sorbonne; I will send it to you when it is published. Before me, Masson-Oursel had made two; he said some implausible things there, and he presented a veritable caricature of Hindu doctrines; if that interests you, I could give you some examples the next time I write you.



\flrightit{Posted on 2012-07-09 by Cologero }

\begin{center}* * *\end{center}

\begin{footnotesize}\begin{sffamily}



\texttt{Avery on 2012-07-09 at 08:27 said: }

I wonder what Guenon thought when the World-Teacher Krishnamurti turned out to be a genius teacher… mired in powerless existentialism.

http://www.kinfonet.org/krishnamurti/excerpts/11/parts


\hfill

\texttt{Cologero on 2012-07-09 at 12:16 said: }

I don't recall that Guenon ever wrote anything about the mature Krishnamurti. Perhaps, because Krishnamurti did not achieve any influence until years later (than the 1920s). I doubt, in any case, that Guenon would have approved of Krishnamurti due to his rejection of traditional forms.


\hfill


\end{sffamily}\end{footnotesize}
